% Copyright (c) 2026 Geovana Neves. Licensed under CC BY 4.0 (https://creativecommons.org/licenses/by/4.0/). See LICENSE-AND-AUTHORSHIP.md.
%
% overview/text/01-idea-and-spectrum.tex
%
% Parts 1 and 2 of guide 1: the solver-agnostic workspace and the spectrum of
% fidelities.

\section{One workspace, many solvers}

\begin{frame}{The problem a workspace solves}
\begin{columns}[T]
\begin{column}{0.48\textwidth}
\begin{block}{Without it}
\begin{itemize}\small
  \item Each solver has its own input deck, its own folder habits and its
        own way of saying that a run ended.
  \item The study lives in scripts that only their writer can read.
  \item Changing solver means starting the bookkeeping again.
\end{itemize}
\end{block}
\end{column}
\begin{column}{0.48\textwidth}
\begin{block}{With it}
\begin{itemize}\small
  \item A study is \textbf{one folder}, the workspace, with the same tree
        for every solver.
  \item A study is \textbf{one table}, the run matrix: one row per
        simulation.
  \item The same commands plan it, run it, follow it and collect it.
\end{itemize}
\end{block}
\end{column}
\end{columns}
\begin{takeaway}You learn the workflow once. The solver changes what is
inside three input files, not how the campaign is run.\end{takeaway}
\end{frame}

\begin{frame}[fragile]{The solver is one line}
\begin{columns}[T]
\begin{column}{0.44\textwidth}
\begin{lstlisting}[style=cfdbtoml]
# init.toml, in the workspace root
schema_version = 1
solver = "su2"
\end{lstlisting}
\begin{itemize}\small
  \item \code{cfd-workspace init} writes it, once.
  \item Every command reads it first.
  \item A workspace keeps its solver for life.
\end{itemize}
\end{column}
\begin{column}{0.52\textwidth}
\relax{\cfdbtablesize
\begin{tabularx}{\linewidth}{@{}P{0.28\linewidth}R@{}}
\toprule
\textbf{Solver name} & \textbf{What runs} \\
\midrule
\code{su2} & SU2, by its executable or its Python wrapper \\
\code{openfoam} & OpenFOAM (ESI), \code{simpleFoam} \\
\code{pyfr} & PyFR, flux reconstruction \\
\code{xfoil} & XFOIL, airfoil polars \\
\code{qprop} & QPROP, an isolated propeller \\
\code{mock} & nothing: a stand-in to learn with \\
\code{flightstream} & FlightStream, handed to pyflightstream \\
\bottomrule
\end{tabularx}}
\end{column}
\end{columns}
\end{frame}

\begin{frame}{What stays the same, what changes}
\relax{\cfdbbodysize
\begin{tabularx}{\linewidth}{@{}P{0.16\linewidth}P{0.4\linewidth}R@{}}
\toprule
 & \textbf{The same for every solver} & \textbf{The solver's own} \\
\midrule
The tree & \code{inputs/}, \code{sims/}, \code{post/}, \code{archive/},
  \code{runs.json} & the files inside one point folder \\
The matrix & its columns and what they mean & a few flight-condition keys \\
The commands & every command and option of \code{cfd-matrix} and
  \code{cfd-workspace} & nothing \\
The inputs & one setup, one reference and one post file per row & the keys
  of the setup file, the mesh format \\
The result & status, verdict, one file set per polar & the coefficient names,
  what ``converged'' rests on \\
\bottomrule
\end{tabularx}}
\vspace{2mm}
\begin{takeaway}Guide 2 teaches the left column. Guide 3 teaches the right
one, solver by solver.\end{takeaway}
\end{frame}

\section{A spectrum of fidelities}

\begin{frame}{More than one wind tunnel}
\begin{itemize}
  \item A laboratory runs more than one wind tunnel: a small one answers many
        cheap questions, a large one answers a few expensive ones.
  \item Solvers are the same. Each resolves a part of the physics and models
        or ignores the rest.
  \item \pkgname{} keeps them behind one workflow, so the choice of fidelity
        is a choice of workspace and not a change of habits.
\end{itemize}
\begin{center}
\begin{tikzpicture}[
  s/.style={draw=hanblue,fill=hanblue!8,rounded corners=2pt,minimum height=11mm,
            text width=18mm,align=center,font=\scriptsize,inner sep=1.5pt},
  d/.style={s,dashed}]
  \node[s] (x) at (0,0) {\textbf{XFOIL}\\2D panel and boundary layer};
  \node[s] (q) at (2.3,0) {\textbf{QPROP}\\blade element};
  \node[d] (f) at (4.6,0) {\textbf{FlightStream}\\3D panel method};
  \node[s] (s) at (6.9,0) {\textbf{SU2}\\Euler, RANS};
  \node[s] (o) at (9.2,0) {\textbf{OpenFOAM}\\RANS};
  \node[s] (p) at (11.5,0) {\textbf{PyFR}\\high order, time resolved};
  \draw[-{Stealth[length=6pt]},thick,SoftGray] (-1.0,-0.95) -- (12.5,-0.95)
    node[midway,below,font=\scriptsize,text=SoftGray]
    {more of the flow resolved, more cost per point};
\end{tikzpicture}
\end{center}
{\scriptsize\color{SoftGray}Dashed: run through pyflightstream, in its own
workspace format.\par}
\end{frame}

\begin{frame}{What each solver answers}
\relax{\cfdbtablesize
\begin{tabularx}{\linewidth}{@{}P{0.13\linewidth}P{0.27\linewidth}P{0.27\linewidth}R@{}}
\toprule
\textbf{Solver} & \textbf{It models} & \textbf{One point is} & \textbf{It
returns} \\
\midrule
XFOIL [5] & a 2D airfoil: panels coupled to an integral boundary layer & one
  angle of attack & \code{CL}, \code{CD}, \code{CDp}, \code{CM} \\
QPROP [6] & an isolated propeller in axial flow, from blade sections and
  their 2D polars & one speed and rotation rate & thrust, torque, power,
  \code{CT}, \code{CP}, efficiency \\
SU2 [7] & compressible flow on a volume mesh, Euler or RANS & one Mach number
  and attitude & force and moment coefficients \\
OpenFOAM [8] & steady incompressible RANS, \code{simpleFoam}, on a 2D
  extruded mesh & one speed and angle of attack & \code{Cd}, \code{Cl},
  \code{CmPitch} \\
PyFR [9] & 2D compressible Navier-Stokes, high order, marched in time & one
  Mach, Reynolds number and angle & time-averaged \code{Cl}, \code{Cd} \\
\bottomrule
\end{tabularx}}
\vspace{1mm}
{\small\color{SoftGray}Every library solver runs \textbf{steady} points
today. Guide 3 states the limits of each solver.\par}
\end{frame}

\begin{frame}{The mock solver}
\begin{columns}[T]
\begin{column}{0.5\textwidth}
\begin{itemize}
  \item \code{mock} is a solver that computes nothing and returns fixed
        numbers.
  \item It needs no installation, no mesh and no licence.
  \item Its setup file can ask for a converged point, a point at the
        iteration limit or a failed one.
\end{itemize}
\end{column}
\begin{column}{0.46\textwidth}
\begin{block}{Use it to}
\begin{itemize}\small
  \item learn the commands (Part 4 of this guide);
  \item try a matrix layout before the meshes exist;
  \item see what a failed point looks like in the records, safely.
\end{itemize}
\end{block}
\end{column}
\end{columns}
\begin{takeaway}Start on the mock. Move to your solver when the workflow is
no longer the thing you are learning.\end{takeaway}
\end{frame}

\begin{frame}{A result is evidence, not validation}
\begin{itemize}
  \item \pkgname{} records what was run, with which inputs, and what the
        solver reported. It keeps the hashes of the inputs beside each
        result.
  \item It classifies each point: \code{converged}, \code{iteration_limit}
        or \code{failed}. What ``converged'' rests on is the solver's own
        evidence, stated in guide 3.
  \item It does \textbf{not} judge the mesh, the turbulence model or the
        physics. A converged point on a poor mesh is still a converged
        point.
\end{itemize}
\begin{cfdbevidence}[What you can rely on]
The same matrix row gives the same input deck, byte for byte, and the record
names every file the point read.
\end{cfdbevidence}
\end{frame}
